% ===== PRÉAMBULE CANONIQUE — à coller VERBATIM en tête de CHAQUE .tex =====
\documentclass[11pt,a4paper]{article}
\usepackage[utf8]{inputenc}\usepackage[T1]{fontenc}\usepackage{lmodern}
\usepackage[french]{babel}
\usepackage{amsmath,amssymb,mathtools}\usepackage{icomma}
\usepackage[version=4]{mhchem}          % \ce{...} : équations & formules
\usepackage{siunitx}\sisetup{locale=FR} % \qty{}{}, \si{}, \ang{}
\usepackage{chemfig}                    % \chemfig{...} : structures développées
\usepackage{xcolor}\usepackage{colortbl}
\usepackage{tikz}\usepackage{pgfplots}\pgfplotsset{compat=1.18}
\usepackage[most]{tcolorbox}\usepackage{enumitem}\usepackage{array,booktabs}
\usepackage[a4paper,margin=2.2cm]{geometry}
\usetikzlibrary{arrows.meta,calc,angles,quotes,positioning,patterns,backgrounds,shapes.geometric}
\definecolor{primary}{RGB}{222,49,99}\definecolor{primarydark}{RGB}{150,22,55}
\definecolor{defcol}{RGB}{0,114,178}\definecolor{methcol}{RGB}{52,92,148}
\definecolor{excol}{RGB}{0,135,90}\definecolor{attncol}{RGB}{200,40,55}
\tcbset{boxbase/.style={enhanced,breakable,boxrule=0.9pt,arc=2.5pt,
  left=3mm,right=3mm,top=2.3mm,bottom=2.3mm,fonttitle=\bfseries,coltitle=white}}
% --- boîtes TUTORIEL (noms EXACTS, ne pas renommer) ---
\newtcolorbox{cle}[1][]{boxbase,colback=defcol!6,colframe=defcol,
  title={\textsc{À retenir}\ifx\relax#1\relax\relax\else\textmd{ : \itshape #1}\fi}}
\newtcolorbox{methode}[1][]{boxbase,colback=methcol!7,colframe=methcol,sharp corners,
  title={\textsc{Méthode}\ifx\relax#1\relax\relax\else\textmd{ --- \itshape #1}\fi}}
\newtcolorbox{exemplebox}[1][]{boxbase,colback=excol!5,colframe=excol,
  title={\textsc{Exemple}\ifx\relax#1\relax\relax\else\textmd{ : \itshape #1}\fi}}
\newtcolorbox{piege}[1][]{boxbase,colback=attncol!6,colframe=attncol,
  title={\textsc{Piège}\ifx\relax#1\relax\relax\else\textmd{ : \itshape #1}\fi}}
\newtcolorbox{remarque}[1][]{boxbase,colback=black!4,colframe=black!45,
  title={\textsc{Remarque}\ifx\relax#1\relax\relax\else\textmd{ : \itshape #1}\fi}}
% --- boîtes EXERCICE / CORRIGÉ (noms EXACTS, auto-comptées) ---
\newtcolorbox[auto counter]{exo}[1][]{boxbase,colback=primary!4,colframe=primary,
  title={\textsc{Exercice}~\thetcbcounter\ifx\relax#1\relax\relax\else\textmd{ --- \itshape #1}\fi}}
\newtcolorbox[auto counter]{corrige}[1][]{boxbase,colback=excol!4,colframe=excol,
  title={\textsc{Corrigé}~\thetcbcounter\ifx\relax#1\relax\relax\else\textmd{ --- \itshape #1}\fi}}
\newcommand{\R}{\mathbb{R}}\newcommand{\N}{\mathbb{N}}\newcommand{\Z}{\mathbb{Z}}\newcommand{\ds}{\displaystyle}
% (un \usepackage{fancyhdr}+en-tête est possible ; il est ignoré côté web)
% =========================================================================
\begin{document}

\begin{exo}[pression totale = somme des partielles]
On applique $p_{\text{tot}}=\sum_i p_i$.
\begin{enumerate}[label=\alph*)]
  \item Un mélange a $p_A=\qty{0.3}{atm}$ et $p_B=\qty{0.5}{atm}$ : calcule $p_{\text{tot}}$.
  \item $p_{\text{tot}}=\qty{2.0}{atm}$, $p_A=\qty{1.2}{atm}$ : calcule $p_B$.
  \item Trois gaz de pressions partielles $\qty{0.2}{atm}$, $\qty{0.3}{atm}$, $\qty{0.5}{atm}$ :
        calcule $p_{\text{tot}}$.
\end{enumerate}
\end{exo}

\begin{exo}[fraction molaire]
On rappelle $\chi_A=\dfrac{n_A}{n_{\text{tot}}}$.
\begin{enumerate}[label=\alph*)]
  \item Un mélange contient $n_A=\qty{2}{\mole}$ et $n_B=\qty{3}{\mole}$ : calcule $\chi_A$ et $\chi_B$.
  \item $n_A=\qty{0.25}{\mole}$ dans un mélange de $n_{\text{tot}}=\qty{1.0}{\mole}$ : calcule $\chi_A$.
  \item Vérifie que $\chi_A+\chi_B=1$ dans le cas a).
\end{enumerate}
\end{exo}

\begin{exo}[pression partielle à partir de la fraction molaire]
On applique $p_A=\chi_A\times p_{\text{tot}}$.
\begin{enumerate}[label=\alph*)]
  \item $\chi_A=0{,}4$, $p_{\text{tot}}=\qty{2.0}{atm}$ : calcule $p_A$.
  \item Air : $\chi(\ce{O2})=0{,}20$, $p_{\text{tot}}=\qty{1.0}{atm}$ : calcule $p(\ce{O2})$.
  \item $\chi_B=0{,}75$, $p_{\text{tot}}=\qty{4.0}{atm}$ : calcule $p_B$.
\end{enumerate}
\end{exo}

\begin{exo}[fraction volumique = fraction molaire]
L'air contient (en volume) $\approx 80\,\%$ de \ce{N2} et $20\,\%$ de \ce{O2}. On prend
$p_{\text{tot}}=\qty{1.0}{atm}$.
\begin{enumerate}[label=\alph*)]
  \item Donne la fraction molaire $\chi(\ce{O2})$.
  \item Donne la fraction molaire $\chi(\ce{N2})$.
  \item Calcule la pression partielle $p(\ce{N2})$.
\end{enumerate}
\end{exo}

\end{document}
